\chapter{What the fluid has and what equation (1) keeps}

This chapter defines the set of every property of a fluid and the set of every force a fluid can
feel, names the members equation (1) keeps, and proves two more ways in which what it keeps is not
what the fluid has. It rests on equations (1) and (2) of Fefferman's
statement~\cite{src:Fefferman}, IAPWS R6-95 (2018)~\cite{src:IAPWS-R6-95-2018}, the formulation
for the thermodynamic properties of water, and Caupin and Herbert's review of cavitation in
water~\cite{src:Caupin-and-Herbert-2006}.

\section{The two sets}

\textbf{Definition.} $P_{\mathrm{all}}$ is the set of all properties of a fluid, known and not yet
known. $F_{\mathrm{all}}$ is the set of all forces a fluid can experience, known and not yet known.
Both sets are open in the sense that no list of their members is claimed complete. A member enters
a list when it is measured, and its absence from a list says only that it has not been measured
yet.

Equation (1), with its density set to one, keeps
\[
P_{(1)} = \{\nu\},
\qquad
F_{(1)} = \{-\nabla p,\ \nu\Delta u,\ f\},
\]
under the constraint $\nabla\cdot u = 0$. The term $(u\cdot\nabla)u$ is the rate of change of
momentum carried with the flow and not a force.

Every member of $P_{\mathrm{all}}\setminus P_{(1)}$ or $F_{\mathrm{all}}\setminus F_{(1)}$ that is
shown present by a measurement shows that equation (1) is a model on part of the two sets. The two sets being
open, no finite list can show the reverse: that equation (1) keeps everything. The previous chapters
exhibit members of both differences, and the workbook collects every one of them with the source
that shows it present, as terms to put back.

\section{Pressure at infinite speed}

\textbf{Proposition 12.} Let $(u,p)$ be smooth and solve (1) and (2) on $\R^3$. At each instant,
\[
\Delta p = \nabla\cdot f - \sum_{i,j}\partial_i\partial_j(u_iu_j).
\]
Among solutions of this equation that decay at infinity, the one $p$ is unique, and its value at
any point depends on $u$ and $f$ at every point of $\R^3$ at the same instant.

\emph{Proof.} Take the divergence of (1). By (2), $\nabla\cdot\partial_t u = \partial_t\theta = 0$ and
$\nabla\cdot(\nu\Delta u) = \nu\Delta\theta = 0$. The remaining term is
$\partial_i(u_j\partial_j u_i) = \partial_i\partial_j(u_iu_j) - \partial_i(u_i\,\partial_j u_j)$,
and the second part is zero by (2). A harmonic function on $\R^3$ that decays at infinity is zero,
which gives uniqueness, and the decaying solution is the convolution of the source with
$-1/(4\pi|x|)$, which is nonzero at every point. \hfill\qedsymbol

A change in the motion anywhere changes the pressure everywhere at the same instant. In water at
$20\,^\circ$C and atmospheric pressure, pressure travels at the speed of sound, 1482.35 meters per
second, evaluated from IAPWS R6-95 (2018) by a program that reproduces its
verification table to every printed digit. The speed of sound is a member of
$P_{\mathrm{all}}\setminus P_{(1)}$, and condition (2) sets it to infinity. The fundamental
solution of the linear problem shows it directly: a force switched on at one instant gives a
pressure proportional to $\delta(t)$ at every point of space~\cite{src:Ladyzhenskaya-1969}.

\section{Pressure with no level}

\textbf{Proposition 13.} If $(u,p)$ solves (1) and (2), then $(u,\,p+c)$ solves them too, for every function $c$
of time alone.

\emph{Proof.} Pressure enters (1) only through $\nabla p$, and $\nabla c = 0$. \hfill\qedsymbol

Equation (1) fixes no level of pressure, and conditions (6) and (7) of Fefferman's statement fix
none either: neither bounds $p$ at infinity. A real liquid has a level at which it breaks. Liquid
water can be held under tension, at negative pressure, until vapor bubbles form, and Caupin and
Herbert report that classical nucleation theory puts that point at $-168$ MPa at $20\,^\circ$C for
the volume and time they state, that water in mineral inclusions is estimated to reach about
$-140$ MPa, and that most other methods reach between a few and a few tens of megapascals below
zero. Whether the liquid breaks depends on the level of the pressure. Equation (1) cannot pose the
question, because it holds no level. Ladyzhenskaya~\cite{src:Ladyzhenskaya-1969} replaces the
requirement that the pressure be non-negative by a bound on its integrals over surfaces, and writes
that those integrals ``only have physical meaning for areas $\Sigma$ whose sizes are not less than a
certain positive number (stipulated by the limits of accuracy of measurement and by the discreteness
of the liquid medium).''

\section{What this establishes and what it does not}

\textbf{Established.} Equation (1) keeps one property and three forces of the two open sets. It
carries pressure at infinite speed, where water carries it at the speed of sound, and it holds no
level of pressure, where water breaks at a level.

\textbf{Not established.} The sets are open, and nothing here bounds what they hold beyond the
members measured. That keeping a member changes whether solutions break down is not shown for any
member in this chapter.
